\documentclass[header={Lecture 9 Exam Study Notes}]{study-note}

\begin{document}

\StudyTitleBlock{09}{Lecture 9: Exam Study Notes}{Fredholm theory and Fredholm alternative}

\vspace{0.8em}

\begin{focusbox}
This lecture is short in scope but central in logic: it explains why operators of the form $I-K$ with $K$ compact are so important, and how solvability reduces to orthogonality conditions involving the adjoint.
\end{focusbox}

\tableofcontents

\section{Setting And Main Question}

\begin{focusbox}
The central equation is
\[
(I-K)u=f,
\]
where $K\in\mathcal L(H)$ is compact and $H$ is a Hilbert space. The lecture explains why this equation behaves much more rigidly than a general linear equation in infinite dimensions.
\end{focusbox}

\begin{defbox}{Homogeneous and inhomogeneous problems}
Associated with
\[
(I-K)u=f
\]
are:
\begin{itemize}
\item the \textbf{homogeneous equation} $(I-K)u=0$;
\item the \textbf{inhomogeneous equation} $(I-K)u=f$.
\end{itemize}
The Fredholm alternative says that the inhomogeneous equation is solvable exactly when $f$ satisfies orthogonality conditions against solutions of the adjoint homogeneous equation.
\end{defbox}

\section{Core Statement}

\begin{thmbox}{Theorem 13.3: Fredholm alternative}
\thmFredholmAlternative
\end{thmbox}

\begin{prooflevelbox}{STATEMENT}
Know the statement precisely: finite-dimensional kernels, closed range, dimension equality, and the dichotomy (unique solution for all $f$ or nontrivial homogeneous solutions). Not the full proof.
\end{prooflevelbox}

\begin{focusbox}
This is the exact exam slogan:
\[
\operatorname{Ran}(I-K)=\ker(I-K^*)^{\perp}.
\]
So solvability is equivalent to orthogonality to all solutions of the adjoint homogeneous equation.
\end{focusbox}

\section{Why Operators Of The Form I Minus K Matter}

\begin{focusbox}
Compact operators are not invertible perturbations in general, but the operator
\[
I-K
\]
behaves in many ways like a finite-dimensional matrix. This is why Fredholm theory feels like an infinite-dimensional version of linear algebra.
\end{focusbox}

\begin{warningbox}{Finite-dimensional analogy}
In finite dimensions, injective and surjective are equivalent for linear maps. Fredholm theory recovers a substitute of this equivalence for $I-K$ when $K$ is compact.
\end{warningbox}

\begin{exbox}{Why compactness is the crucial hypothesis}
For a general bounded operator $T$, injectivity does \emph{not} imply surjectivity and the range need not be closed. For compact perturbations of the identity, the defect is finite-dimensional and the range has a controlled form.
\end{exbox}

\section{Structural Facts Behind The Alternative}

\begin{thmbox}{Riesz--Schauder structure for $I-K$}
If $K\in\mathcal L(H)$ is compact, then:
\begin{enumerate}
\item $\ker(I-K)$ is finite-dimensional;
\item $\ker(I-K^*)$ is finite-dimensional;
\item $\operatorname{Ran}(I-K)$ is closed;
\item
\[
\operatorname{Ran}(I-K)=\ker(I-K^*)^{\perp}.
\]
\end{enumerate}
These facts are what make the Fredholm alternative possible.
\end{thmbox}

\paragraph{Interpretation.}
The obstruction to solving $(I-K)u=f$ is finite-dimensional: only finitely many orthogonality conditions must be checked.

\begin{thmbox}{Uniqueness versus solvability}
For compact $K$, the following implications are central:
\begin{itemize}
\item if $\ker(I-K)=\{0\}$, then the solution of $(I-K)u=f$ is unique whenever it exists;
\item if $\ker(I-K^*)=\{0\}$, then every $f\in H$ satisfies the orthogonality condition, so the equation is solvable for every $f\in H$;
\item in the Fredholm situation, injectivity and surjectivity are tied together much more closely than for general bounded operators.
\end{itemize}
\end{thmbox}

\section{Adjoint Formulation Of Solvability}

\begin{defbox}{Adjoint homogeneous equation}
The adjoint equation is
\[
(I-K^*)v=0.
\]
Its solution space is $\ker(I-K^*)$.
\end{defbox}

\begin{thmbox}{Orthogonality criterion}
The equation
\[
(I-K)u=f
\]
is solvable if and only if
\[
(f,v)_H=0 \qquad \forall v\in \ker(I-K^*).
\]
So every adjoint homogeneous solution gives one linear compatibility condition on the right-hand side.
\end{thmbox}

\begin{exbox}{How to use this in practice}
To decide whether $(I-K)u=f$ is solvable:
\begin{enumerate}
\item solve or understand the adjoint homogeneous problem $(I-K^*)v=0$;
\item list a basis of $\ker(I-K^*)$;
\item check the finitely many conditions $(f,v)_H=0$ on that basis.
\end{enumerate}
If all conditions hold, then a solution exists.
\end{exbox}

\section{Role Of The Homogeneous Equation}

\begin{focusbox}
The homogeneous equation
\[
(I-K)u=0
\]
asks whether $1$ is an eigenvalue of $K$.
\end{focusbox}

\begin{thmbox}{Equivalent viewpoint}
The equation $(I-K)u=0$ has a nonzero solution if and only if there exists $u\neq 0$ such that
\[
Ku=u.
\]
Thus the failure of uniqueness for $(I-K)u=f$ is exactly the presence of the eigenvalue $1$ for $K$.
\end{thmbox}

\paragraph{Exam interpretation.}
Lecture 9 is already preparing Lecture 10: spectral questions about $\lambda I-K$ are closely connected with eigenvalues of compact operators.

\section{What To Remember About The Proof Structure}

\begin{exbox}{Key ingredients}
The proof uses:
\begin{itemize}
\item compactness of $K$;
\item the adjoint operator $K^*$;
\item the fact that compactness controls the structure of kernels and ranges;
\item the identity $\operatorname{Ran}(I-K)=\ker(I-K^*)^{\perp}$;
\item orthogonality of the right-hand side to the kernel of the adjoint.
\end{itemize}
You are usually expected to know the theorem, its meaning, and the logical route of the proof, not every technical step in full detail.
\end{exbox}

\begin{warningbox}{Where compactness enters}
One of the main ingredients from the previous lecture is that compact operators send weakly convergent sequences to strongly convergent sequences. This is used in the background to show that $I-K$ has much better range properties than a general bounded operator.
\end{warningbox}

\section{Consequences To State Cleanly In An Exam}

\begin{thmbox}{Standard consequences}
Let $K\in\mathcal L(H)$ be compact.
\begin{enumerate}
\item Either $(I-K)u=f$ has a unique solution for every $f\in H$, or the homogeneous equation $(I-K)u=0$ has a nontrivial solution.
\item If $\ker(I-K^*)\neq\{0\}$, then solvability requires orthogonality conditions on $f$.
\item If $\ker(I-K)=\{0\}$, then there is no nontrivial homogeneous solution, hence solutions of the inhomogeneous problem are unique.
\item If both kernels are trivial, then $(I-K)$ is bijective.
\end{enumerate}
\end{thmbox}

\begin{exbox}{Finite-dimensional picture to keep in mind}
Think of a square matrix $A$. The system $Ax=b$ is solvable exactly when $b$ is orthogonal to all vectors in $\ker(A^*)$, and uniqueness fails exactly when $\ker(A)\neq\{0\}$. Fredholm theory says that $I-K$ with $K$ compact behaves in the same spirit.
\end{exbox}

\section{Exam-Style Exercises}

\begin{exbox}{Exercise 9.1: Fredholm compatibility condition}
Let $H$ be a Hilbert space, let $K\in\mathcal L(H)$ be compact, and suppose that
\[
\ker(I-K^*)=\operatorname{span}\{v_1,\dots,v_m\}.
\]
State necessary and sufficient conditions on $f\in H$ for the equation
\[
(I-K)u=f
\]
to be solvable.

\paragraph{Solution pattern.}
Use
\[
\operatorname{Ran}(I-K)=\ker(I-K^*)^{\perp}.
\]
Therefore solvability is equivalent to
\[
(f,v_j)_H=0 \qquad j=1,\dots,m.
\]
\end{exbox}

\begin{exbox}{Exercise 9.2: Uniqueness versus kernel}
Let $K\in\mathcal L(H)$ be compact. Prove that if $\ker(I-K)=\{0\}$, then $(I-K)u=f$ has a unique solution for every $f\in H$.

\paragraph{Solution pattern.}
Fredholm theory says injectivity and surjectivity are linked for $I-K$. Trivial kernel gives uniqueness, and the Fredholm alternative gives solvability for every $f$. Then $(I-K)$ is bijective.
\end{exbox}

\begin{exbox}{Exercise 9.3: Nonuniqueness of solutions}
Assume $(I-K)u=f$ has one solution $u_0$ and $\ker(I-K)\neq\{0\}$. Describe all solutions.

\paragraph{Solution pattern.}
Show that every solution has the form
\[
u_0+w,
\qquad w\in\ker(I-K).
\]
Subtract two solutions to prove the converse inclusion.
\end{exbox}

\begin{exbox}{Exercise 9.4: A finite-rank Fredholm example}
Let $H=L^2(0,1)$ and define
\[
(Kf)(x)=\left(\int_0^1 f(t)\,dt\right)\cdot 1.
\]
\begin{enumerate}
\item Show that $K$ is compact.
\item Compute $K^*$.
\item Solve the homogeneous equation $(I-K)u=0$.
\item Determine for which $f\in L^2(0,1)$ the equation $(I-K)u=f$ is solvable.
\end{enumerate}

\paragraph{Solution pattern.}
$K$ has rank one, so it is compact. It is selfadjoint because it is defined by orthogonal projection onto constants up to scaling. The homogeneous equation gives $u=Ku$, so $u$ must be constant. The solvability condition is
\[
f\perp 1,
\]
that is,
\[
\int_0^1 f(t)\,dt=0.
\]
\end{exbox}

\begin{exbox}{Exercise 9.5: Relation to eigenvalues}
Show that $(I-K)u=0$ has a nonzero solution if and only if $1$ is an eigenvalue of $K$.

\paragraph{Solution pattern.}
Simply rewrite
\[
(I-K)u=0
\iff
Ku=u.
\]
\end{exbox}

\begin{exbox}{Exercise 9.6: Why compactness matters}
Give an example of a bounded operator $T$ on an infinite-dimensional Hilbert space such that $I-T$ is injective but not surjective.

\paragraph{Solution pattern.}
Take the unilateral right shift $S$ on $\ell^2$. Then $I-S$ is injective because $x=Sx$ forces $x=0$, but it is not surjective: for example, $e_1$ has no preimage in $\ell^2$ because solving $(I-S)x=e_1$ gives the non-square-summable sequence $(1,1,1,\dots)$.
\end{exbox}

\section{What To Memorize Precisely}

\begin{focusbox}
For the exam, know:
\begin{enumerate}
\item the statement of the Fredholm alternative;
\item the structural facts that $\ker(I-K)$ and $\ker(I-K^*)$ are finite-dimensional and that $\operatorname{Ran}(I-K)$ is closed;
\item the solvability condition involving $\ker(I-K^*)$;
\item the identity $\operatorname{Ran}(I-K)=\ker(I-K^*)^{\perp}$;
\item why nonuniqueness is equivalent to the homogeneous equation having a nontrivial solution;
\item that $(I-K)u=0$ means $1$ is an eigenvalue of $K$;
\item why the special form $I-K$ with $K$ compact is essential;
\item the finite-dimensional analogy.
\end{enumerate}
\end{focusbox}

\end{document}
